\documentclass[11pt]{article}
\usepackage{mystyle}

\title{FIT1058 --- Chapter 3: Proofs\\ \large Exercise Solutions}
\author{Alston}
\date{Semester 2, 2026}

\begin{document}
\maketitle

% ======================================================================
\begin{exercise}{1}
    Recall the Collatz function
    \[
        \Collatz(x) =
        \begin{cases}
            3x+1 & \text{if } x \text{ is odd} \\
            \dfrac{x}{2} & \text{if } x \text{ is even}
        \end{cases}
    \]
    Prove by cases that, for all positive integers $n$,
    \[
        \Collatz^{(2)}(n) \le \tfrac{3}{2}n + 1.
    \]
\end{exercise}
\begin{proof}
    Proof by cases.
    \begin{align*}
        \begin{cases}
            x \text{ is odd, }  x=2n+1\\
            x \text{ is even, } x=2m
        \end{cases}
    \end{align*}
    \begin{enumerate}[label=(\Roman*)]
        \item $x$ is odd,
    \end{enumerate}
\end{proof}

% ======================================================================
\begin{exercise}{2}
    A set of strings $L$ is called \emph{hereditary} if it has the
    following property: for every nonempty string $x$ in $L$, there is
    a character in $x$ which can be deleted from $x$ to give another
    string in $L$.

    Prove by contradiction that every nonempty hereditary set of
    strings contains the empty string.
\end{exercise}
\begin{solution}

\end{solution}

% ======================================================================
\begin{exercise}{5}
    Let $E$ be a finite set and let $\I$ be a collection of subsets of
    $E$ that satisfies the following three properties:
    \begin{description}
        \item[(I1)] $\emptyset \in \I$.
        \item[(I2)] Whenever $X \subseteq Y$ and $Y \in \I$, we also
            have $X \in \I$. This is called the \emph{hereditary
            property}.
        \item[(I3)] Whenever $X, Y \in \I$ and $|X| < |Y|$, there
            exists $y \in Y \setminus X$ such that $X \cup \{y\} \in
            \I$. In other words, given two different-sized sets in
            $\I$, you can always enlarge the smaller set, using some
            new member from the larger set, to get another set in
            $\I$. This is called the \emph{independence augmentation
            property}.
    \end{description}
    The sets in $\I$ are called \emph{independent}.

    For example, if $E = \{a,b,c\}$, then each of the following
    collections satisfies all three properties (I1)--(I3):
    \begin{itemize}
        \item $\I = \{\emptyset, \{a\}, \{b\}, \{a,b\}\}$
        \item $\I = \{\emptyset, \{a\}, \{b\}, \{c\}\}$
        \item $\I = \{\emptyset, \{a\}, \{b\}, \{c\}, \{a,c\},\{b,c\}\}$
        \item $\I = \{\emptyset, \{a\}, \{b\}, \{c\}, \{a,b\},\{a,c\},\{b,c\}\}$
        \item $\I = \{\emptyset\}$
    \end{itemize}
    By way of counterexample, using the same $E$, each of the
    following collections fails to satisfy all three properties
    (I1)--(I3):
    \begin{itemize}
        \item $\I = \{\{a\}, \{b\}, \{a,b\}\}$. This fails (I1) and
            (I2).
        \item $\I = \{\emptyset, \{a\}, \{a,b\}\}$. This fails (I2).
        \item $\I = \{\emptyset, \{a\}, \{b\}, \{c\}, \{a,b\}\}$. This
            satisfies (I1) and (I2) but fails (I3).
        \item $\I = \emptyset$. This fails (I1).
    \end{itemize}

    \begin{enumerate}[label=(\alph*)]
        \item For each counterexample above that fails (I2) and/or
            (I3), give specific sets $X$ and $Y$ for which the
            property fails.

        \item Prove that, for a social network, the set of cliques
            satisfies (I1) and (I2) but not, in general, (I3).

        \item Prove by contradiction that all maximal independent
            sets have the same size.

            Suppose that each member $e \in E$ has a nonnegative real
            weight $w(e)$, so $w: E \to \R_0^+$. We use this weight
            function to also give weights to subsets of $E$, by
            defining the weight of a set to be the sum of the weights
            of its elements: for all $X \subseteq E$, define
            \[
                w(X) := \sum_{e \in X} w(e).
            \]
            Suppose we seek an independent set of maximum weight.
            Consider the following \emph{greedy algorithm}, so called
            because it always makes the choice that gives biggest
            immediate gain:

            \begin{algorithm}[H]
                \DontPrintSemicolon
                \KwIn{a finite set $E$, a function $w: E \to \R_0^+$, a collection $\I$ of subsets of $E$}
                \KwOut{an independent set $X$}
                $X \gets \emptyset$\;
                $i \gets 1$\;
                \While{$\exists\, y \in E \setminus X$ such that $X \cup \{y\} \in \I$}{
                    choose $y$ with largest weight $w(y)$ among all such candidates\;
                    $e_i \gets y$\;
                    $X \gets X \cup \{e_i\}$\;
                    $i \gets i + 1$\;
                }
                \Return $X$\;
                \caption{Greedy algorithm for maximum weight independent set}
            \end{algorithm}

        \item Construct a simple social network, and give a
            nonnegative real weight to each person, such that the
            greedy algorithm, with $\I$ being the set of all cliques,
            does not find the maximum weight clique.

            Now suppose $\I$ satisfies (I1)--(I3). The greedy
            algorithm chooses elements of $E$ in a specific order,
            represented as $e_1, e_2, \dots, e_r$.

        \item Prove by contradiction that the weights of the chosen
            elements are decreasing, i.e.,
            \[
                w(e_1) \ge w(e_2) \ge \cdots \ge w(e_r).
            \]
            (``Decreasing'' allows the possibility that some weights
            stay the same; we don't require strictly decreasing.)

        \item Prove that the greedy algorithm finds a maximum weight
            independent set.
    \end{enumerate}
\end{exercise}
\begin{solution}[(a)]

\end{solution}
\begin{solution}[(b)]

\end{solution}
\begin{solution}[(c)]

\end{solution}
\begin{solution}[(d)]

\end{solution}
\begin{solution}[(e)]

\end{solution}
\begin{solution}[(f)]

\end{solution}
\FloatBarrier
% ======================================================================
\begin{exercise}{10}
    Consider the following algorithm:

    \begin{algorithm}[H]
        \DontPrintSemicolon
        \KwIn{a finite set $E$, a function $w: E \to \R_0^+$, a collection $\I$ of subsets of $E$}
        \KwOut{an independent set $X$}
        $X \gets \emptyset$\;
        $i \gets 1$\;
        \While{$\exists\, y \in E \setminus X$ such that $X \cup \{y\} \in \I$}{
            choose $y$ with largest weight $w(y)$ among all such candidates\;
            $e_i \gets y$\;
            $X \gets X \cup \{e_i\}$\;
            $i \gets i + 1$\;
        }
        \Return $X$\;
        \caption{Greedy algorithm for maximum weight independent set}
    \end{algorithm}
    \begin{enumerate}[label=(\alph*)]
        \item For each pair $i,j$, how many times is the action
            inside the innermost loop performed? Express this in
            terms of $i$ or $j$ or both.
        \item For each $i$, how many times is the action inside the
            innermost loop performed? Write this as both a sum of
            some number of terms, and (using a theorem in this
            chapter) a simple algebraic expression, both in terms of
            $i$.
        \item Now write a sum of some number of terms, in terms of
            $n$, for the total number of times the action is
            performed.
        \item Work out this expression for $n = 1,2,3,4,5$, and for as
            many further values of $n$ as you like.
        \item Explore how these numbers behave.
        \item Conjecture a simple algebraic expression, in terms of
            $n$, for the total number of times the action is
            performed.
        \item Prove, by induction on $n$, that your expression is
            correct.
    \end{enumerate}
\end{exercise}
\begin{solution}

\end{solution}

% ======================================================================
\begin{exercise}{11}
    The $n$-th harmonic number $H_n$ is defined by
    \[
        H_n = 1 + \frac{1}{2} + \frac{1}{3} + \cdots + \frac{1}{n}.
    \]
    In this exercise we prove by induction that $H_n \ge \ln(n+1)$.
    (It follows that the harmonic numbers increase without bound, even
    though $H_n$ grows more and more slowly as $n$ increases.)
    \begin{enumerate}[label=(\roman*)]
        \item \textbf{Inductive basis:} prove that $H_1 \ge \ln(1+1)$.
        \item Let $n \ge 1$. Assume $H_n \ge \ln(n+1)$ (inductive
            hypothesis). Consider $H_{n+1}$; write it recursively
            using $H_n$. Use the inductive hypothesis to obtain
            $H_{n+1} \ge \dots$, then use the fact that $\ln(1+x) \le
            x$, and a property of logarithms, to show
            $H_{n+1} \ge \ln(n+2)$.
        \item What can you now conclude, and why?
    \end{enumerate}
    \textbf{Advanced afterthoughts.}
    \begin{itemize}
        \item The inequality implies $H_n \ge \ln n$, since
            $\ln(n+1) \ge \ln n$. It is instructive to try to prove
            this directly by induction --- you will probably run into
            a snag, illustrating that induction sometimes requires
            proving something stronger than what you set out to
            prove.
        \item Would your proof work for logarithms to other bases,
            apart from $e$? Where does the proof use base $e$?
        \item It is known that $H_n \le (\ln n) + 1$. Can you prove
            this?
    \end{itemize}
\end{exercise}
\begin{solution}

\end{solution}

% ======================================================================
\begin{exercise}{12}
    (Challenge) Let $E$ be a finite set, and let $\mathcal{S}$ be a
    set of subsets of $E$. We say $\mathcal{S}$ is \emph{linear under
    $\triangle$} if, for every two sets $X, Y \in \mathcal{S}$, we
    have $X \triangle Y \in \mathcal{S}$.

    The \emph{span} of $\mathcal{S}$, written $\mathrm{span}(\mathcal{S})$,
    is the set of all symmetric differences of any number of members
    of $\mathcal{S}$:
    \[
        \mathrm{span}(\mathcal{S}) := \{X_1 \triangle X_2 \triangle
        \cdots \triangle X_k : X_1,\dots,X_k \in \mathcal{S},\ k \in
        \N\}.
    \]
    By convention, the symmetric difference of zero sets is
    $\emptyset$, so $\mathrm{span}(\mathcal{S})$ always contains
    $\emptyset$.

    $\mathcal{S}$ is \emph{dependent under $\triangle$} if, for some
    distinct $X_1,\dots,X_k \in \mathcal{S}$ (i.e.\ $X_i \ne X_j$ for
    $i \ne j$), $X_1 \triangle \cdots \triangle X_k = \emptyset$. If
    $\mathcal{S}$ is not dependent under $\triangle$, it is
    \emph{independent under $\triangle$}.

    $\mathcal{S}$ is a \emph{circuit under $\triangle$} if it is
    dependent under $\triangle$ and also minimal with respect to that
    property. $\mathcal{S}$ is a \emph{base under $\triangle$} if it
    is independent under $\triangle$ and also maximal with respect to
    that property.

    \begin{enumerate}[label=(\alph*)]
        \item Let $\mathcal{S}$ be linear under $\triangle$. Prove, by
            induction on $k$, that for all $k$ and all $X_1,\dots,X_k
            \in \mathcal{S}$, $X_1 \triangle \cdots \triangle X_k \in
            \mathcal{S}$.
        \item Prove that, for every set $\mathcal{S}$ of subsets of
            $E$, $\mathcal{S}$ is linear under $\triangle$ if and only
            if $\mathcal{S} = \mathrm{span}(\mathcal{S})$.
        \item Prove by induction on $k$ that for all $k$ and all
            $X_1,\dots,X_k \in \mathcal{S}$, either
            $X_1\triangle\cdots\triangle X_k = \emptyset$, or there
            exist distinct $Y_1,\dots,Y_n \in \mathcal{S}$ with $n \le
            k$ such that $X_1\triangle\cdots\triangle X_k =
            Y_1\triangle\cdots\triangle Y_n$.
        \item Prove that $\mathcal{S}$ is dependent under $\triangle$
            if and only if there exists $X \in \mathcal{S}$ such that
            $X \in \mathrm{span}(\mathcal{S}\setminus\{X\})$.
        \item Prove that if $\mathcal{S}$ is a minimal dependent set
            under $\triangle$, then for every $X \in \mathcal{S}$ we
            have $X \in \mathrm{span}(\mathcal{S}\setminus\{X\})$.
        \item Prove that every minimal spanning set under $\triangle$
            is independent under $\triangle$.
        \item Prove that every minimal spanning set under $\triangle$
            is a maximal independent set under $\triangle$.
        \item Prove that every maximal independent set under
            $\triangle$ is also a spanning set under $\triangle$.
        \item Prove that every maximal independent set under
            $\triangle$ is also a minimal spanning set under
            $\triangle$.
        \item Prove that $\mathcal{S}$ is a minimal spanning set under
            $\triangle$ if and only if it is a maximal independent set
            under $\triangle$.
        \item Prove that, if $\mathcal{S}$ is independent under
            $\triangle$ then $|\mathrm{span}(\mathcal{S})| =
            2^{|\mathcal{S}|}$.
        \item Prove the independence augmentation property: if
            $\mathcal{S}$ and $\mathcal{T}$ are independent under
            $\triangle$ and $|\mathcal{T}| = |\mathcal{S}|+1$, then
            there exists $X \in \mathcal{T}\setminus\mathcal{S}$ such
            that $\mathcal{S}\cup\{X\}$ is independent under
            $\triangle$.
        \item Prove that all bases under $\triangle$ have the same
            size.
        \item Do all minimal dependent sets under $\triangle$ have the
            same size? Give a proof of your claim.
        \item Prove that, for every base $\mathcal{B}$ of $\mathcal{S}$
            under $\triangle$, and every $X \in \mathcal{S}$, there is
            a unique way to write $X$ as a symmetric difference of
            members of $\mathcal{B}$, using each element of
            $\mathcal{B}$ at most once.
    \end{enumerate}
    For each previous part (a)--(o), state what type of proof you have
    used.

    Finally: design a simple algorithm for the following problem ---
    given a collection $\mathcal{S}$ of subsets of $E$ with a
    nonnegative real weight for each member, output a subset
    $\mathcal{X}$ of $\mathcal{S}$ that is independent under
    $\triangle$ such that the sum of the weights of all members of
    $\mathcal{X}$ is maximum. Does this algorithm always find the best
    possible solution? Justify your answer, using a previous exercise.
\end{exercise}
\begin{solution}[(a)]

\end{solution}
\begin{solution}[(b)]

\end{solution}
\begin{solution}[(c)]

\end{solution}
\begin{solution}[(d)]

\end{solution}
\begin{solution}[(e)]

\end{solution}
\begin{solution}[(f)]

\end{solution}
\begin{solution}[(g)]

\end{solution}
\begin{solution}[(h)]

\end{solution}
\begin{solution}[(i)]

\end{solution}
\begin{solution}[(j)]

\end{solution}
\begin{solution}[(k)]

\end{solution}
\begin{solution}[(l)]

\end{solution}
\begin{solution}[(m)]

\end{solution}
\begin{solution}[(n)]

\end{solution}
\begin{solution}[(o)]

\end{solution}
\begin{solution}[Algorithm \& justification]

\end{solution}

% ======================================================================
\begin{exercise}{13}
    Identify the errors in the following ``theorem'' and ``proof'',
    which are both incorrect. The steps in the ``proof'' are numbered
    for convenience.

    \medskip
    \textbf{``Theorem.''} For all $n \in \N$, a circular disc can be
    cut into $2^n$ pieces by $n$ straight lines.

    \textbf{``Proof.''}
    \begin{enumerate}
        \item Consider $n=0$. The disc is not cut at all, so it is
            still in one piece. So the number of pieces is $2^0=1$.
            So the claim is true for $n=0$.
        \item Consider $n=1$. A single line across a disc cuts it
            into two pieces, so the number of pieces is $2^1$. So the
            claim is true for $n=1$.
        \item Now consider $n=2$. If the circle is cut by one line
            into two pieces, then we can cut it again by a line that
            crosses the first line. Such a line must cut right across
            both of the pieces created by the first cut. Therefore
            each of those pieces is divided in two, so the total
            number of pieces is $2\times 2 = 2^2$, so the claim is
            true for $n=2$.
        \item It is clear from the cases considered so far that, once
            the circle is cut by some number of lines, then another
            line can be used to divide every piece into two pieces,
            thereby doubling the number of pieces.
        \item So, if we use $n$ lines, then the repeated doubling
            gives $2^n$ pieces. Therefore the claim is true for all
            $n$. \hfill ``$\square$''
    \end{enumerate}
\end{exercise}
\begin{solution}

\end{solution}

% ======================================================================
\begin{exercise}{14}
    Identify the errors in the following ``theorem'' and ``proof'',
    which are both incorrect.

    \medskip
    \textbf{``Theorem.''} Every rational number is an integer.

    \textbf{``Proof.''}
    \begin{enumerate}
        \item We need to show that the sets $\Z$ and $\Q$ are equal.
        \item To prove that two sets are equal, we can prove the
            appropriate subset and superset relations between the two
            sets.
        \item We first prove the subset relation between these two
            sets, $\Z \subseteq \Q$.
        \item Let $x \in \Z$.
        \item Since $x$ is an integer, we have $x = \dfrac{x}{1}$.
        \item So $x$ is a quotient of integers.
        \item Therefore $x$ is rational, i.e., $x \in \Q$.
        \item So we have proved that $\Z \subseteq \Q$.
        \item Now we need to prove the superset relation, which is
            the reverse of subset.
        \item So we will prove that $\Q \supseteq \Z$.
        \item Let $x \notin \Q$.
        \item Then $x$ cannot be written in the form $\dfrac{p}{q}$
            where $p$ and $q$ are integers.
        \item So it certainly cannot be written in this form with
            $q=1$.
        \item Therefore $x$ is not an integer, i.e., $x \notin \Z$.
        \item So we have proved that $\Q \supseteq \Z$.
        \item It follows from our subset and superset relations
            between the two sets that $\Z = \Q$. \hfill ``$\square$''
    \end{enumerate}
    \begin{quote}
        \emph{``God made the integers; all else is the work of man.''}
        \hfill --- Leopold Kronecker (1823--1891)
    \end{quote}
\end{exercise}
\begin{solution}

\end{solution}

% ======================================================================
\begin{exercise}{15}
    Let $G$ be the cryptosystem obtained from Caesar slide by
    restricting the keyspace to the first half of the alphabet (see
    Exercise 2.15 for definitions relating to the Caesar slide
    cryptosystem). So the keyspace $K_G$ is
    \[
        \{\mathrm{a,b,c,d,e,f,g,h,i,j,k,l,m}\}.
    \]
    Prove that $G$ is not idempotent.
\end{exercise}
\begin{solution}

\end{solution}

\end{document}